\subsection{Principal Component Analysis (PCA)}

What if we have $n$ data points, $p$ features and $n<p$ ?
It means that some fatures might look something completely different, but in 
reality some of them are just a combination of others $\rightarrow$ we 
have \textbf{collinear features}.
Collinear features make harder our job of understanding how each feature influences
the targets : let's assume that $x_1$ and $x_2$ are collinear
\begin{equation}
    y = \beta_0 + \beta_1 x_1 + \beta_2 x_2 +\varepsilon
\end{equation}
due collinearity, we can't see the individual effect of $x_1$ and $x_2$ on $y$.
%\vspace{0.8cm}

\begin{figure}[h]
    \centering
    \includegraphics[width=0.5\textwidth]{dr/pca/dr_matrix_example.png}    
\end{figure}
What we want to do is to summarize the data having $p$ features with the smaller
set ok $k$ artifically derived features; to do this we must balance the clarity of
representation and the loss of information due the compression $\Rightarrow$ 
\textbf{Principal Component Analysis (PCA)} aims to do this using linear 
combinations of the selected variables.

\vspace{0.6cm}
\textcolor{Green}{\textbf{Example}:}\\
Imagine that we want to describe our blood cells sample, and the initial 
features are \textit{diameter} and \textit{thickness}, how can we reduce
the number of features?
\begin{figure}[h]
    \centering    
    \includegraphics[width=0.3\textwidth]{dr/pca/blood_cells.png}
\end{figure}
\newpage

We can think of projecting the points along one of the two dimensions, choosing it
based on how spreaded the data are along it (e.g. so preserving more information).
\begin{center} \textcolor{red}{
    \textbf{In statitics information = variability.}   
} \end{center}
\begin{figure} [h]
    \centering
    \begin{subfigure}[c]{0.25\textwidth}
        \centering
        \includegraphics[width=\linewidth]{dr/pca/dr_blood_ex1.png}
    \end{subfigure}
    \hspace{1cm}
    \begin{subfigure}[c]{0.2\textwidth}
        \centering
        \includegraphics[width=\linewidth]{dr/pca/dr_blood_ex2.png}
    \end{subfigure}
\end{figure}

So we need a way to capture the maximal variability of the data: in other
words to find \textit{along which dimension the data spread the most}.
\begin{figure}[h]
    \centering
    \includegraphics[width=0.3\textwidth]{dr/pca/pca_maxvar_example.png}
\end{figure}

PCA's goal is to project the data along those directions of maximal variability,
ending up with a new reference system for the data points \textcolor{gray}{(it's 
nothing different from vector space rotations, in 2D you visualize it by immaging
that the axis rotate)}.\\
\hfill \\

This objective can be expressed in two equivalent ways:
\begin{description}
    \item[Maximize variance]: maximize the data spread along that
    direction
    \item[Minimize residuals]: minimize the distance of the data points
    from the the dimension
\end{description}
\begin{figure}[h]
    \centering
    \includegraphics[width=0.45\textwidth]{dr/pca/pca_goal_visualzed.png}
\end{figure}

In order to do the PCA:
\begin{enumerate}
    \item get a set of data points $x^{(1)}\dots x^{(n)}\in \real^p$
    \item center the data in $(0,0)$ if needed: $x^{(i)}\leftarrow x^{(i)}-
    \bar{x}$
    \item choose a direction $u\in\real^p\ :\ \|x\|_2=1$
    \item project the data $x^{(i)}$ on $u$ : $z_i = u^Tx^{(i)}\in \real$
\end{enumerate}

Taking the objective of maximizing the variance, we observe that since the data are 
centered (avg = 0), the varince is:
\begin{equation*}
    Var(z) = \frac{1}{n}\sum_{i=1}^n z_i^2 = \frac{1}{n}\sum_{i=1}^n 
    (u^Tx^{(i)})^2 = 0
    \ \rightarrow\ 
    \max_{u\in\real^p}\frac{1}{n}\sum_{i=1}^n (u^Tx^{(i)})^2
\end{equation*}
if we define the \textit{empirical covariance matrix} as
\begin{equation*}
    \Sigma = \frac{1}{n}\sum_{i=1}^n x^{(i)}x^{(i)T} \ \in 
    \real^{p\times p}
\end{equation*}
we get that
\begin{equation*}
    \frac{1}{n}\sum_{i=1}^n (u^Tx^{(i)})^2 = u^T(\frac{1}{n}
    \sum_{i=1}^n x^{(i)}x^{(i)T})u = u^T\Sigma u
    \rightarrow
    \max_{u\in\real^p} u^T\Sigma u
\end{equation*}
The constraint on $u$ can be expressed as $u^Tu=1$ (turning constraint into 
an equation)

By switching to the Lagrangian relaxed problem:
\begin{equation*}
    \mathcal{L}(u,\lambda)=u^T\Sigma u-\lambda(u^Tu-1)
\end{equation*}

and taking the derivative
\begin{equation*}
    \nabla_u (u^T\Sigma u) = 2\Sigma u,\ \nabla_u (u^Tu) = 2u
    \Rightarrow \nabla_u \mathcal(L) = 2\Sigma u-2\lambda u = 0
    \Rightarrow \Sigma u = \lambda u
\end{equation*}

$\Rightarrow$ a candidate direction $u$ must be an \textbf{eigenvector of $\mathbf{\Sigma}$}.\\
So the objective function becomes $u^T\Sigma u=u^T (\lambda u) = \lambda
(u^T u) = \lambda \rightarrow$ by choosing the highest $\lambda$ we get the
highest objective function. So as principal component we take the eigenvector with the highest 
eigenvalue. If we want more dimension, we take also the second, the third ... \\
\hfill \\

We can interpret $\Sigma$ so that, when applied to the data
modifies the shape, it strech them toward their principal component, which 
are the $\Sigma$'s eigenvectors.
\begin{figure}[h]
    \centering
    \includegraphics[width=0.3\textwidth]{dr/pca/covmat_projection.png}
\end{figure}

Two nice properties of PCA are:
\begin{enumerate}
    \item PCA results are rotation-invariant (depend only on the data, not 
    on the starting axis)
    \item PCs are independent/uncorrelated from each other by construction (they are eigenvectors)
\end{enumerate}

\newpage
\subsubsection{Express data with PCA's components}
We can see that the original observations can be re-written using the new
components: given the starting dimensions $\langle x_1\dots x_p \rangle$, we 
have
\begin{equation*}
    PC_j = \phi_{j1}x_1 \dots + \phi_{jp}x_p \Rightarrow z_{ij} = \phi_{1j}x_{i1}
    \dots + \phi_{pj}x_{ip}
\end{equation*}
where $z_{i,j}$ is the projection of $x^{(i)}$ on the $PC_j$ called \textbf{score}
and $\phi_{1j},\dots \phi_{pj}$ are the \textbf{weights} given to each original 
feature to translate it into the new PCA coordinates system. 
\begin{figure}[h]
    \centering
    \includegraphics[width=0.7\textwidth]{dr/pca/pca_projection_example.png}
\end{figure}

\hfill \\
Note: standardization ($x-\mu/\sigma$) address the problem of
high-variance features dominating principal components.
\vspace{0.5cm}


\textbf{Why this stuff is useful other than re-shaping the data?}\\
Usually the data variability depends only on a small subset of all $p$ features, 
so we can use these to have a $k$-low-dimensional/compressed representation of 
the dataset. The choice of $k$ is the trade-off.

\subsubsection{Portion of Variance Explained}
We would like to quantify how much each singular Principal Component (PC)
explains the data variability.
\begin{equation*}
    Var_{\scriptscriptstyle TOT}= \sum_{j=1}^p Var(x_j) = \sum_{j=1}^p \left(\frac{1}{n}\sum_{i=1}^n x_{ij}^2\right)
\end{equation*}
If we name $Var_k$ := variance explained by the single k-th PC and $Var_{exp}$ := the
portion of Variance Explained by the selected first $k$ components
\begin{equation*}
    Var_k = \frac{1}{n}\sum_{i=1}^n z_{ik}^2 \Rightarrow Var_{exp} = \cfrac
    {\sum_k Var_k}{Var_{\scriptscriptstyle TOT}}
\end{equation*}
However, since the PCs are the eigenvalues of $\Sigma$:
\begin{equation*}
    Var_k = \frac{1}{n}\sum_{i=1}^n z_{ik}^2 = u_k^T\Sigma u_k = u_k^T(\lambda_k
    u_k) = \lambda_k
\end{equation*}
\begin{equation*}
    Var_{\scriptscriptstyle TOT} = \sum_{j=1}^p Var(x_j) = tr(\Sigma) = \sum_{j=1}^p \lambda_k
\end{equation*}
So we can define:
\begin{description}
    \item[$\mathbf{EVR_k}$ Explained Variance Ratio of $\mathbf{k}$] = $\cfrac{\lambda_k}
    {\sum_{m=1}p}\lambda_m$
    \item[$\mathbf{CumEV(K)}$ Cumulative Explained Variance] of the $K$ selection of
    component = $\cfrac{\sum_{k=1}^K\lambda_k}{\sum_{m=1}p}\ \lambda_m$   
\end{description}

\newpage
Depending on how much variability we'd like to retain, we can module the number
of dimensions in $K$ based on the \textit{CumEV index}.
This technique can be used also to other purposes, not only to tune the variance
level: it can be used also to set the data noise treshold to remove those components which 
"make too much noise".
\begin{figure}[h]
    \centering
    \includegraphics[width=0.4\textwidth]{dr/pca/cev_visualized.png}
\end{figure}


\subsubsection{How to interpret PCs}
The PC's weights can be used to interpret the meaning of the PCs itself: by loking at them we can 
have a clue about which dimensions go along in the same direction, and from them we can give to the PC some kind of interpetation
\begin{figure}[h]
    \centering
    \includegraphics[width=0.4\textwidth]{dr/pca/interpret_pcs.png}
\end{figure}
\hfill\\

Another useful tool to interpret the PCs' meaning is the \textbf{biplot}: using
PCs as axis, we plot both the points projected on the PCA dimensions and the original
features.
\begin{figure}[h]
    \centering
    \includegraphics[width=0.4\textwidth]{dr/pca/biplot_example.png}
\end{figure} \hfill \\
and we interpret the meaning as following:
\begin{itemize}
    \item \textbf{arrow direction}: along which direction the feature increases
    \item \textbf{arrow length}: how strongly the features is represented by
    these PCs
    \item \textbf{angle between arrows}: relationship between original features
\end{itemize}

